\documentclass[10pt,a4paper]{amsart}
\usepackage[margin=19mm]{geometry}
\usepackage{amsmath,amssymb,mathtools}
\usepackage[scr=rsfs,cal=euler]{mathalfa}
\usepackage{xfrac}
\usepackage[unicode,hidelinks,bookmarks=false]{hyperref}
\newcommand{\ie}{i.e.\ }
\newcommand{\cf}{cf.\ }
\newcommand{\resp}{respectively\ }
\newcommand{\Subsection}{\subsection}
\providecommand{\nobreakdash}{\nobreak\hbox{-}}
\setlength{\emergencystretch}{2em}
\multlinegap0pt
\usepackage{amssymb}
\newtheorem{theorem}{Theorem}
\title{Nonparametric Inference for Cumulative Residual Mathai--Haubold Entropy of order $\alpha$}
\author{Anija C.R, Smitha, S., Sudheesh, K. Kattumannil}
\date{Source: arXiv:2609.08434v1 (CC BY 4.0)}
\begin{document}
\maketitle

\noindent\emph{Source and licence.} Nonparametric Inference for Cumulative Residual Mathai--Haubold Entropy of order $\alpha$. Version arXiv:2609.08434v1 (CC BY 4.0), distributed by arXiv under the Creative Commons Attribution 4.0 International licence, http://creativecommons.org/licenses/by/4.0/, as recorded in the arXiv licence metadata for this exact version and shown on the abstract page https://arxiv.org/abs/2609.08434v1. Full licence notice: \texttt{licenses/arxiv-2609.08434-CC-BY-4.0.txt}.\par\smallskip

\noindent\textit{Editorial scope.} Counted unit: the article's theorem theorem (source0.tex lines 725--730). The article's complete original proof of it is delivered with it (source0.tex lines 731--766, one \texttt{proof} environment of 36 lines). No mathematics is written, repaired or re-derived: the statement and the proof are the article's own lines, in the article's own notation, and no retained line is substituted or re-flowed.  Retained uncounted as the source-local context the proof needs: lines 699--702.  Nothing was omitted from any retained span, so the package carries no omission record.  No retained line contains a citation, so the wrapper carries no bibliography and no bibliography entry.  No retained line contains a figure environment, an image-inclusion command or drawing code, so no image is delivered and the package declares no asset dependency.  Editorial formatting only, in addition: an AMS \texttt{amsart} preamble that restates the article's own declarations and macros used by the retained lines, namely the environments \texttt{theorem} on separate counters, together with the standard packages those lines need.  The retained environments print in this document as Theorem 1 in order of appearance, whereas the article prints the same items with the numbering of its own full document.  The complete native source of the article as distributed in the arXiv e-print for this exact version is bundled unchanged as \texttt{sources/209715/source0.tex}.
\begin{equation}\label{relationship}
\zeta_{\alpha}'(X_t)
=\bigl((2-\alpha)\zeta_{\alpha}(X_t)-m(t)\bigr)\,h(t).
\end{equation}

\begin{theorem}
 Consider a non-negative random variable $X$ with probability density function $f(x)$,
survival function $\bar{F}(x)$, hazard rate $h(x)$, and mean residual life function
$m(x)$. Suppose that $\zeta_{\alpha}(X_t)$ is increasing in $t$. Then
$\zeta_{\alpha}(X_t)$ uniquely determines the survival function $\bar{F}(x)$.
\end{theorem}

 \begin{proof}
Suppose that $F(x)$ and $G(x)$ are the distribution functions of the random variables
$X$ and $Y$, respectively, and assume that
\begin{equation}\label{A CRM F=G}
\zeta_{\alpha}(X_t)=\zeta_{\alpha}(Y_t), \qquad t \ge 0 .
\end{equation}
Differentiating both sides of \eqref{A CRM F=G} with respect to $t$ and using
\eqref{relationship}, we obtain
\begin{equation}\label{hazard}
\bigl((2-\alpha)\zeta_{\alpha}(X_t)-m_{1}(t)\bigr)\,h_{1}(t)
=
\bigl((2-\alpha)\zeta_{\alpha}(Y_t)-m_{2}(t)\bigr)\,h_{2}(t),
\end{equation}
where $h_{1}(t)$ and $h_{2}(t)$ denote the hazard rates, and $m_{1}(t)$ and $m_{2}(t)$
denote the mean residual life functions corresponding to $X$ and $Y$, respectively.

If $h_{1}(t)=h_{2}(t)$ for all $t \ge 0$, then $\bar{F}(t)=\bar{G}(t)$ for all
$t \ge 0$, and hence $F=G$. Therefore, it suffices to show that
$h_{1}(t)=h_{2}(t)$ for all $t \ge 0$.

Suppose, to the contrary, that there exists $t_0 \ge 0$ such that
$h_{1}(t_0)\neq h_{2}(t_0)$. Without loss of generality, assume that
$h_{1}(t_0)>h_{2}(t_0)$. Then, from \eqref{hazard}, we obtain
\[
(2-\alpha)\zeta_{\alpha}(X_{t_0})-m_{1}(t_0)
<
(2-\alpha)\zeta_{\alpha}(Y_{t_0})-m_{2}(t_0),
\]
which implies that
\[
m_{1}(t_0)>m_{2}(t_0).
\]
This contradicts Lemma~3.1, which states that $h_{1}(t_0)>h_{2}(t_0)$ implies
$m_{1}(t_0)<m_{2}(t_0)$. Hence, such a $t_0$ cannot exist, and therefore
$h_{1}(t)=h_{2}(t)$ for all $t \ge 0$. Consequently, $F=G$, which completes the proof.
\end{proof}

\end{document}
